\section*{Phụ lục}
\label{sec:appendix}

% =======================================================================
% A.1 CẤU HÌNH THỰC NGHIỆM CHI TIẾT
% =======================================================================
\subsection*{A.1 Cấu hình Thực nghiệm Chi tiết}
\label{subsec:experiment_config}

\subsubsection*{A.1.1 Thông số Kiến trúc Mô hình}
\label{subsubsec:model_architecture}

Bảng \ref{tab:model_configs} tổng hợp các thông số kiến trúc cho cả hai mô hình được sử dụng trong thực nghiệm:

\begin{table}[H]
\centering
\caption{So sánh chi tiết cấu hình kiến trúc mô hình}
\label{tab:model_configs}
\begin{tabular}{lcc}
\toprule
\textbf{Thông số} & \textbf{Mini-BERT (From-scratch)} & \textbf{BERT-Base (Fine-tuned)} \\
\midrule
Số lớp Transformer (L) & 2 & 12 \\
Kích thước ẩn (H) & 128 & 768 \\
Số attention heads (A) & 2 & 12 \\
Kích thước FFN & 512 & 3072 \\
Kích thước từ vựng & 30,522 & 30,522 \\
Độ dài chuỗi tối đa & 512 & 512 \\
Số loại segment & 2 & 2 \\
Dropout rate & 0.1 & 0.1 \\
\bottomrule
\end{tabular}
\end{table}

\subsubsection*{A.1.2 Hyperparameters}
\label{subsubsec:training_hyperparams}

\paragraph{BERT configuration (bảng trên)}
\begin{table}[H]
\centering
\caption{Hyperparameters cho BERT}
\label{tab:hyperparameters}
\begin{tabular}{lcc}
\toprule
\textbf{Hyperparameter} & \textbf{BERT "from-scratch"} & \textbf{BERT-Base Fine-tuning} \\
\midrule
Learning rate & 5e-5 & 2e-5 \\
Batch size & 4 & 16 \\
Gradient accumulation steps & 1 & 2 \\
Number of epochs & 10 & 3 \\
Warmup steps & 500 & 500 \\
Weight decay & 0.01 & 0.01 \\
Optimizer & AdamW & AdamW \\
Adam $\beta_1$ & 0.9 & 0.9 \\
Adam $\beta_2$ & 0.999 & 0.999 \\
Adam $\epsilon$ & 1e-8 & 1e-8 \\
Max gradient norm & 1.0 & 1.0 \\
Learning rate schedule & Linear decay & Linear decay \\
Mixed precision (FP16) & Yes & Yes \\
\bottomrule
\end{tabular}
\end{table}

\paragraph{Baseline configuration:}
Mô hình baseline (TF-IDF + Logistic Regression) sử dụng các tham số:
\begin{itemize}
    \item Max features: 5,000
    \item N-gram range: (1, 2)
    \item Stop words: English stopwords
    \item Solver: liblinear
    \item Max iterations: 1,000
    \item Random state: 42
\end{itemize}

% =======================================================================
% A.2 MÔI TRƯỜNG HARDWARE/SOFTWARE
% =======================================================================
\subsection*{A.2 Môi trường}
\label{subsec:environment}

\subsubsection*{A.2.1 Cấu hình Hardware}
\label{subsubsec:hardware_config}

\begin{table}[H]
\centering
\caption{Chi tiết của môi trường thực nghiệm}
\label{tab:hardware_specs}
\begin{tabular}{ll}
\toprule
\textbf{Component} & \textbf{Specification} \\
\midrule
Operating System & Google Colab \\
GPU & NVIDIA Tesla T4 (16GB VRAM) \\
CUDA Version & 11.7 / 12.2 \\
\bottomrule
\end{tabular}
\end{table}

\subsubsection*{A.2.2 Dependencies}
\label{subsubsec:software_deps}

Toàn bộ dependencies được quản lý thông qua file \texttt{requirements.txt}:
\begin{minted}{bash}
torch>=2.0.0
transformers>=4.25.0
matplotlib>=3.5.0
seaborn>=0.11.0
tqdm>=4.64.0
numpy>=1.21.0
scikit-learn>=1.2.0
datasets>=2.8.0
fsspec>=2023.1.0
\end{minted}

% =======================================================================
% A.3 KẾT QUẢ THỰC NGHIỆM MỞ RỘNG
% =======================================================================
\subsection*{A.3 Kết quả Thực nghiệm Mở rộng}
\label{subsec:extended_results}

\subsubsection*{A.3.1 Metrics Chi tiết}
\label{subsubsec:detailed_metrics}

\begin{table}[H]
\centering
\caption{Kết quả chi tiết trên tập validation IMDb (1,000 samples)}
\label{tab:detailed_results}
\begin{tabular}{lccccc}
\toprule
\textbf{Model} & \textbf{Accuracy} & \textbf{Precision} & \textbf{Recall} & \textbf{F1-Score} & \textbf{Training Time} \\
\midrule
TF-IDF + LR & 0.855 & 0.86 & 0.86 & 0.85 & 2 minutes \\
BERT Fine-tuned & \textbf{0.894} & \textbf{0.89} & \textbf{0.90} & \textbf{0.89} & 4.5 minutes \\
\midrule
\textbf{Improvement} & \textbf{+3.9\%} & \textbf{+3.0\%} & \textbf{+4.0\%} & \textbf{+4.0\%} & - \\
\bottomrule
\end{tabular}
\end{table}

\begin{table}[H]
\centering
\caption{Hiệu năng theo từng class (Positive/Negative)}
\label{tab:per_class_results}
\begin{tabular}{lcccccc}
\toprule
\textbf{Model} & \textbf{Class} & \textbf{Precision} & \textbf{Recall} & \textbf{F1-Score} & \textbf{Support} \\
\midrule
\multirow{2}{*}{TF-IDF + LR} & Negative & 0.87 & 0.84 & 0.85 & 506 \\
& Positive & 0.84 & 0.87 & 0.86 & 494 \\
\midrule
\multirow{2}{*}{BERT Fine-tuned} & Negative & 0.90 & 0.89 & 0.89 & 506 \\
& Positive & 0.89 & 0.90 & 0.89 & 494 \\
\bottomrule
\end{tabular}
\end{table}

\subsubsection*{A.3.2 Confusion Matrices}
\label{subsubsec:confusion_matrices}

\begin{table}[H]
\centering
\caption{Confusion Matrix - BERT Fine-tuned Model}
\label{tab:confusion_matrix_bert}
\begin{tabular}{cc|cc}
\toprule
& & \multicolumn{2}{c}{\textbf{Predicted}} \\
& & \textbf{Negative} & \textbf{Positive} \\
\midrule
\multirow{2}{*}{\rotatebox{90}{\textbf{Actual}}} & \textbf{Negative} & 447 & 53 \\
& \textbf{Positive} & 53 & 447 \\
\bottomrule
\end{tabular}
\end{table}

\begin{table}[H]
\centering
\caption{Confusion Matrix - TF-IDF + Logistic Regression Baseline}
\label{tab:confusion_matrix_baseline}
\begin{tabular}{cc|cc}
\toprule
& & \multicolumn{2}{c}{\textbf{Predicted}} \\
& & \textbf{Negative} & \textbf{Positive} \\
\midrule
\multirow{2}{*}{\rotatebox{90}{\textbf{Actual}}} & \textbf{Negative} & 425 & 81 \\
& \textbf{Positive} & 64 & 430 \\
\bottomrule
\end{tabular}
\end{table}

\subsubsection*{A.3.3 Learning Curves}
\label{subsubsec:learning_curves}

\begin{table}[H]
\centering
\caption{Qúa trình train của BERT Fine-tuning}
\label{tab:training_progression}
\begin{tabular}{cccc}
\toprule
\textbf{Epoch} & \textbf{Training Loss} & \textbf{Validation Loss} & \textbf{Validation Accuracy} \\
\midrule
1 & 0.1204 & 0.5563 & 0.8770 \\
2 & 0.0996 & 0.5123 & 0.8940 \\
3 & 0.1045 & 0.6461 & 0.8930 \\
\bottomrule
\end{tabular}
\end{table}

\begin{table}[H]
\centering
\caption{Qúa trình train BERT "from-scratch":}
\label{tab:minibert_training}
\begin{tabular}{cccc}
\toprule
\textbf{Epoch} & \textbf{Average Loss} \\
\midrule
1 & 11.0717 \\
2 & 10.9688 \\
5 & 10.8573 \\
10 & 10.7116 \\
\bottomrule
\end{tabular}
\end{table}

% =======================================================================
% A.4 CODE IMPLEMENTATION  
% =======================================================================
\subsection*{A.4 Code Implementation}
\label{subsec:code_implementation}

\subsubsection*{A.4.1 Các thành phần của BERT}
\label{subsubsec:bert_architecture}

\paragraph{Multi-Head Self-Attention:}
\begin{minted}{python}
class MultiHeadSelfAttention(nn.Module):
    def __init__(self, config: BertConfig):
        super().__init__()
        # Kiểm tra hidden_size chia hết cho num_heads
        self.num_attention_heads = config.num_attention_heads
        self.attention_head_size = int(config.hidden_size / config.num_attention_heads)
        self.all_head_size = self.num_attention_heads * self.attention_head_size

        # Linear projections cho Q, K, V
        # W^Q, W^K, W^V trong paper
        self.query = nn.Linear(config.hidden_size, self.all_head_size)
        self.key = nn.Linear(config.hidden_size, self.all_head_size)
        self.value = nn.Linear(config.hidden_size, self.all_head_size)

        # Dropout cho attention probabilities
        self.dropout = nn.Dropout(config.attention_probs_dropout_prob)

    def transpose_for_scores(self, x: torch.Tensor) -> torch.Tensor:
        new_x_shape = x.size()[:-1] + (self.num_attention_heads, self.attention_head_size)
        x = x.view(*new_x_shape)
        return x.permute(0, 2, 1, 3)

    def forward(self, hidden_states: torch.Tensor, 
                attention_mask: Optional[torch.Tensor] = None) -> Tuple[torch.Tensor, torch.Tensor]:
        # Tính Q, K, V projections
        mixed_query_layer = self.query(hidden_states)
        mixed_key_layer = self.key(hidden_states)
        mixed_value_layer = self.value(hidden_states)

        # Reshape cho multi-head: (batch, heads, seq_len, head_size)
        query_layer = self.transpose_for_scores(mixed_query_layer)
        key_layer = self.transpose_for_scores(mixed_key_layer)
        value_layer = self.transpose_for_scores(mixed_value_layer)

        # Tính attention scores: QK^T
        # shape: (batch, heads, seq_len, seq_len)
        attention_scores = torch.matmul(query_layer, key_layer.transpose(-1, -2))
        
        # Scale bởi √d_k để tránh softmax saturation khi d_k lớn
        # Đây là "Scaled" trong "Scaled Dot-Product Attention"
        attention_scores = attention_scores / math.sqrt(self.attention_head_size)

        if attention_mask is not None:
            # Áp dụng mask: -10000 cho positions cần mask
            # Sau softmax sẽ thành ~0
            attention_scores = attention_scores + attention_mask

        # Normalize scores thành probabilities
        attention_probs = nn.Softmax(dim=-1)(attention_scores)
        
        # Dropout cho regularization
        attention_probs = self.dropout(attention_probs)

        # Tính weighted sum: attention_probs @ V
        context_layer = torch.matmul(attention_probs, value_layer)

        # Reshape về (batch, seq_len, all_head_size)
        context_layer = context_layer.permute(0, 2, 1, 3).contiguous()
        new_context_layer_shape = context_layer.size()[:-2] + (self.all_head_size,)
        context_layer = context_layer.view(*new_context_layer_shape)

        return context_layer, attention_probs
\end{minted}

\paragraph{BERT Embeddings Layer:}
\begin{minted}{python}
class BertEmbeddings(nn.Module):
    def __init__(self, config: BertConfig):
        super().__init__()
        # Token embeddings với padding_idx=0 ([PAD] token)
        self.word_embeddings = nn.Embedding(
            config.vocab_size, 
            config.hidden_size, 
            padding_idx=0
        )
        
        # Position embeddings - BERT học position thay vì dùng sinusoidal
        self.position_embeddings = nn.Embedding(
            config.max_position_embeddings, 
            config.hidden_size
        )
        
        # Segment embeddings để phân biệt câu A/B, Cần thiết cho NSP task
        self.token_type_embeddings = nn.Embedding(
            config.type_vocab_size, 
            config.hidden_size
        )
        
        # LayerNorm với epsilon nhỏ tránh chia 0
        self.LayerNorm = nn.LayerNorm(config.hidden_size, eps=1e-12)
        
        # Dropout để regularization
        self.dropout = nn.Dropout(config.hidden_dropout_prob)

    def forward(self, input_ids: torch.Tensor, 
                token_type_ids: Optional[torch.Tensor] = None) -> torch.Tensor:
        seq_length = input_ids.size(1)
        
        # Tạo position IDs từ 0 đến seq_length-1
        position_ids = torch.arange(seq_length, dtype=torch.long, device=input_ids.device)
        position_ids = position_ids.unsqueeze(0).expand_as(input_ids)

        # Default segment 0 nếu không có token_type_ids
        if token_type_ids is None:
            token_type_ids = torch.zeros_like(input_ids)

        # Lấy 3 loại embeddings
        word_embeddings = self.word_embeddings(input_ids)
        position_embeddings = self.position_embeddings(position_ids)
        token_type_embeddings = self.token_type_embeddings(token_type_ids)

        # Cộng tổng 3 embeddings (Figure 2 trong paper)
        embeddings = word_embeddings + position_embeddings + token_type_embeddings

        # LayerNorm và Dropout
        embeddings = self.LayerNorm(embeddings)
        embeddings = self.dropout(embeddings)
        
        return embeddings
\end{minted}

\subsubsection*{A.4.2 Training Pipeline}
\label{subsubsec:training_pipeline}

\paragraph{Dataset tự tạo:}
\begin{minted}{python}
small_corpus = [
    ("The man went to the store.", "He bought a gallon of milk."),
    ("She reads a book.", "The book is about dragons."),
    ("BERT is a powerful model.", "It was developed by Google."),
    ("An apple a day keeps the doctor away.", "The cat sat on the mat."),
    ("The dog is cute.", "The sky is blue."),
    ("A transformer is a type of model.", "This model uses an attention mechanism."),
    ("This is a language model.", "The model can understand text."),
    ("They built a new AI model.", "The model has many parameters."),
    ("This system uses a predictive model.", "The model forecasts future sales."),
    ("The new model is very efficient.", "It runs faster than the old one."),
    ("GPT-3 is a large language model.", "The model was trained by OpenAI."),
    ("This machine learning model is complex.", "The model requires a lot of data."),
    # --- Dữ liệu bổ sung để dạy về câu liên quan (NSP) ---
    ("I have a new car.", "The car is red."),
    ("The student studied for the exam.", "He passed with a high score."),
    ("We went to the restaurant last night.", "The food there was delicious."),
    ("She wrote a long letter.", "The letter was for her best friend."),
    ("The library is quiet.", "People are reading books."),
    # --- Dữ liệu bổ sung để dạy về câu không liên quan (NSP) ---
    ("The sun is shining today.", "My favorite color is green."),
    ("He drinks coffee every morning.", "History is an interesting subject."),
    ("The train arrived on time.", "Elephants are very large mammals."),
    ("She likes to play the piano.", "The ocean is very deep."),
    ("My computer is very fast.", "Birds can fly high in the sky.")
]
\end{minted}

\paragraph{Chuẩn bị dataset pre-traing cho BERT "from-scratch":}
\begin{minted}{python}
class PretrainingDataset(Dataset):
    def __init__(self, corpus, tokenizer, max_len=64):
        self.corpus = corpus
        self.tokenizer = tokenizer
        self.max_len = max_len

    def __len__(self):
        return len(self.corpus)

    def __getitem__(self, idx):
        # === Chuẩn bị cho NSP ===
        sent_a, sent_b = self.corpus[idx]
        is_next = True
        
        # 50% xác suất chọn random sentence
        if random.random() < 0.5:
            rand_idx = random.randint(0, len(self.corpus) - 1)
            sent_b = self.corpus[rand_idx][1]
            is_next = False

        # Tokenize với format [CLS] A [SEP] B [SEP]
        tokenized = self.tokenizer(
            sent_a, sent_b, 
            truncation=True, 
            max_length=self.max_len, 
            padding="max_length"
        )
        
        input_ids = torch.tensor(tokenized['input_ids'])
        token_type_ids = torch.tensor(tokenized['token_type_ids'])
        attention_mask = torch.tensor(tokenized['attention_mask'])

        # === Chuẩn bị cho MLM ===
        labels = input_ids.clone()
        
        # Tạo probability matrix 15% cho masking
        probability_matrix = torch.full(labels.shape, 0.15)
        
        # Không mask special tokens ([CLS], [SEP], [PAD])
        special_tokens_mask = torch.tensor(
            self.tokenizer.get_special_tokens_mask(labels.tolist(), already_has_special_tokens=True), 
            dtype=torch.bool
        )
        probability_matrix.masked_fill_(special_tokens_mask, value=0.0)
        
        # Chọn positions để mask dựa trên xác suất
        masked_indices = torch.bernoulli(probability_matrix).bool()
        
        # Set label -100 cho non-masked positions (ignored trong loss)
        labels[~masked_indices] = -100

        # Áp dụng chiến lược 80-10-10
        # 80% masked positions -> [MASK]
        indices_replaced = torch.bernoulli(torch.full(labels.shape, 0.8)).bool() & masked_indices
        input_ids[indices_replaced] = self.tokenizer.mask_token_id

        # 10% masked positions -> random token  
        indices_random = torch.bernoulli(torch.full(labels.shape, 0.5)).bool() & masked_indices & ~indices_replaced
        random_words = torch.randint(len(self.tokenizer), labels.shape, dtype=torch.long)
        input_ids[indices_random] = random_words[indices_random]
        
        # 10% còn lại giữ nguyên

        return {
            "input_ids": input_ids,
            "token_type_ids": token_type_ids,
            "attention_mask": attention_mask,
            "masked_lm_labels": labels,
            "next_sentence_label": torch.tensor(1 if is_next else 0, dtype=torch.long)
        }
\end{minted}


\paragraph{Training BERT "from-scratch":}
\begin{minted}{python}
def train_loop(model, dataloader, optimizer, epochs, device):
    model.train()
    print("Bắt đầu quá trình huấn luyện BERT...")
    
    for epoch in range(epochs):
        total_loss = 0
        progress_bar = tqdm(dataloader, desc=f"Epoch {epoch + 1}/{epochs}")
        
        for batch in progress_bar:
            optimizer.zero_grad()
            
            # Move batch lên device
            inputs = {k: v.to(device) for k, v in batch.items()}
            
            # Forward pass
            outputs = model(**inputs)
            loss = outputs['loss']
            
            if loss is not None:
                # Backward pass
                loss.backward()
                
                # Gradient clipping (optional nhưng recommended)
                torch.nn.utils.clip_grad_norm_(model.parameters(), 1.0)
                
                # Update weights
                optimizer.step()
                
                total_loss += loss.item()
                
            # Update progress bar
            progress_bar.set_postfix({'loss': loss.item() if loss else 'N/A'})

        avg_loss = total_loss / len(dataloader)
        print(f"Epoch {epoch + 1} hoàn tất. Loss trung bình: {avg_loss:.4f}")
        
    print("Huấn luyện hoàn tất!")
\end{minted}

% =======================================================================
% A.5 PHÂN TÍCH TOÁN HỌC
% =======================================================================
\subsection*{A.5 Phân tích toán}
\label{subsec:mathematical_analysis}

\subsubsection*{A.5.1 Attention}
\label{subsubsec:attention_math}

\paragraph{Scaled Dot-Product Attention:}

Công thức cơ bản của self-attention trong BERT:

\begin{equation}
\text{Self-Attention}(Q, K, V) = \text{softmax}\left(\frac{QK^T}{\sqrt{d_k}}\right)V
\label{eq:scaled_attention}
\end{equation}

Trong đó:
\begin{align}
Q &= XW^Q \in \mathbb{R}^{n \times d_k} \quad \text{(Queries)} \\
K &= XW^K \in \mathbb{R}^{n \times d_k} \quad \text{(Keys)} \\
V &= XW^V \in \mathbb{R}^{n \times d_v} \quad \text{(Values)} \\
X &\in \mathbb{R}^{n \times d_{model}} \quad \text{(Input representations)}
\end{align}

\paragraph{Multi-Head Attention:}
\begin{equation}
\text{MultiHead}(Q, K, V) = \text{Concat}(\text{head}_1, ..., \text{head}_h)W^O
\label{eq:multihead_full}
\end{equation}

Với mỗi head $i$:
\begin{equation}
\text{head}_i = \text{Self-Attention}(QW_i^Q, KW_i^K, VW_i^V)
\label{eq:head_i_full}
\end{equation}

Các ma trận chiếu tham số:
\begin{align}
W_i^Q &\in \mathbb{R}^{d_{model} \times d_k}, \quad d_k = d_{model}/h \\
W_i^K &\in \mathbb{R}^{d_{model} \times d_k} \\
W_i^V &\in \mathbb{R}^{d_{model} \times d_v}, \quad d_v = d_{model}/h \\
W^O &\in \mathbb{R}^{hd_v \times d_{model}}
\end{align}

\paragraph{Độ phức tạp:}

self-attention:
\begin{align}
\text{Time Complexity} &= O(n^2 \cdot d + n \cdot d^2) \\
\text{Space Complexity} &= O(n^2 + n \cdot d)
\end{align}


\subsubsection*{A.5.2 Hàm Loss}
\label{subsubsec:loss_functions}

\paragraph{Masked Language Model Loss:}
\begin{equation}
\mathcal{L}_{\text{MLM}} = -\sum_{i \in \mathcal{M}} \log P(x_i | \mathbf{x}_{\backslash \mathcal{M}})
\label{eq:mlm_loss}
\end{equation}

Trong đó $\mathcal{M}$ là tập các vị trí bị mask, và $\mathbf{x}_{\backslash \mathcal{M}}$ là sequence với các tokens bị mask.

\paragraph{Next Sentence Prediction Loss:}
\begin{equation}
\mathcal{L}_{\text{NSP}} = -\sum_{i=1}^{N} y_i \log P(\text{IsNext} | C, A_i, B_i)
\label{eq:nsp_loss}
\end{equation}

Với $y_i \in \{0, 1\}$ là label cho sentence pair $(A_i, B_i)$.

\paragraph{Combined Pre-training Loss:}
\begin{equation}
\mathcal{L}_{\text{total}} = \mathcal{L}_{\text{MLM}} + \mathcal{L}_{\text{NSP}}
\label{eq:combined_loss}
\end{equation}

\subsubsection*{A.5.3 Optimization}
\label{subsubsec:optimization}

\paragraph{AdamW Optimizer:}

BERT sử dụng AdamW optimizer với decoupled weight decay:

\begin{align}
m_t &= \beta_1 m_{t-1} + (1 - \beta_1) g_t \\
v_t &= \beta_2 v_{t-1} + (1 - \beta_2) g_t^2 \\
\hat{m}_t &= \frac{m_t}{1 - \beta_1^t} \\
\hat{v}_t &= \frac{v_t}{1 - \beta_2^t} \\
\theta_t &= \theta_{t-1} - \alpha \left(\frac{\hat{m}_t}{\sqrt{\hat{v}_t} + \epsilon} + \lambda \theta_{t-1}\right)
\end{align}

Với: $\beta_1 = 0.9$, $\beta_2 = 0.999$, $\epsilon = 10^{-8}$, $\lambda = 0.01$.

\paragraph{Cập nhật Learning Rate:}
Linear warmup và decay:
\begin{equation}
lr(t) = \begin{cases}
lr_{max} \cdot \frac{t}{t_{warmup}} & \text{if } t \leq t_{warmup} \\
lr_{max} \cdot \frac{t_{total} - t}{t_{total} - t_{warmup}} & \text{if } t > t_{warmup}
\end{cases}
\label{eq:lr_schedule}
\end{equation}

% =======================================================================
% A.6 DATASET VÀ PREPROCESSING
% =======================================================================
\subsection*{A.6 Dataset}
\label{subsec:dataset_preprocessing}

\subsubsection*{A.6.1 Tổng quan IMDb dataset}
\label{subsubsec:imdb_statistics}

\begin{table}[H]
\centering
\caption{Thống kê chi tiết bộ dữ liệu IMDb Movie Reviews}
\label{tab:imdb_statistics}
\begin{tabular}{lc}
\toprule
\textbf{Thống kê} & \textbf{Giá trị} \\
\midrule
Tổng số reviews & 50,000 \\
Training set (official) & 25,000 \\
Test set (official) & 25,000 \\
\midrule
\textbf{Subset được sử dụng:} & \\
Training samples & 4,000 \\
Validation samples & 1,000 \\
\midrule
Positive/Negative distribution & 50\%/50\% \\
\midrule
Vocabulary size & ~89,000 \\
After WordPiece tokenization & ~120,000\\
\bottomrule
\end{tabular}
\end{table}

\subsubsection*{A.6.2 Preprocessing}
\label{subsubsec:preprocessing_steps}

\paragraph{Tokenizer cho BERT fine-tune:}
\begin{minted}{python}
tokenizer = BertTokenizer.from_pretrained('bert-base-uncased')

def preprocess_function(examples):
    return tokenizer(
        examples['text'],
        truncation=True,
        padding='max_length',
        max_length=256,  # Optimized cho GPU memory
        return_tensors='pt'
    )
\end{minted}

% =======================================================================
% A.7 PHÂN TÍCH LỖI
% =======================================================================
\subsection*{A.7 Phân tích Lỗi}
\label{subsec:error_analysis}

\begin{table}[H]
\centering
\caption{Ví dụ các trường hợp sarcasm/irony dự đoán sai}
\label{tab:sarcasm_examples}
\begin{tabular}{p{8cm}cc}
\toprule
\textbf{Review Text (excerpt)} & \textbf{True} & \textbf{Predicted} \\
\midrule
"This movie is just truly awful, the eye-candy that plays Ben just can make up for everything else that is wrong with this movie." & Negative & Positive \\
\midrule
"Thanks to this fungal film I do NOT want my Maypo, and that's saying something." & Negative & Positive \\
\midrule
"What a brilliant piece of film making. So original and thought-provoking. NOT!" & Negative & Positive \\
\bottomrule
\end{tabular}
\end{table}

\begin{table}[H]
\centering
\caption{Ví dụ complex negation patterns}
\label{tab:negation_examples}
\begin{tabular}{p{8cm}cc}
\toprule
\textbf{Text} & \textbf{True} & \textbf{Predicted} \\
\midrule
"I previously thought that this film was the lamest thing ever... I was never more wrong about anything in my life." & Positive & Negative \\
\midrule
"You can't not love this movie if you don't hate good cinematography." & Positive & Negative \\
\midrule
"It's not that I didn't enjoy it, but rather that it wasn't uninteresting." & Positive & Negative \\
\bottomrule
\end{tabular}
\end{table}

% =======================================================================
% A.8 TRỰC QUAN HÓA
% =======================================================================
\subsection*{A.8 Trực quan hóa}
\label{subsec:visualization}

\subsubsection*{A.8.1 t-SNE cho embeddings}
\label{subsubsec:tsne_embeddings}

\begin{table}[H]
\centering
\caption{Cấu hình t-SNE cho visualization}
\label{tab:tsne_params}
\begin{tabular}{ll}
\toprule
\textbf{Parameter} & \textbf{Value} \\
\midrule
n\_components & 2 \\
perplexity & 30 \\
learning\_rate & 200 \\
n\_iter & 1000 \\
random\_state & 42 \\
metric & euclidean \\
early\_exaggeration & 12.0 \\
\bottomrule
\end{tabular}
\end{table}

\paragraph{Quy trình trích xuất BERT embedding và visualize:}
\begin{minted}{python}
# Tải mô hình tốt nhất đã được lưu
best_model_path = trainer.state.best_model_checkpoint
print(f"Đang tải mô hình tốt nhất từ: {best_model_path}")
feature_extractor = BertModel.from_pretrained(best_model_path).to(device)
feature_extractor.eval()

# Lấy một tập con nhỏ hơn để trực quan hóa
sample_dataset = tokenized_val_dataset.shuffle(seed=RANDOM_STATE).select(range(500))

def get_embeddings(batch):
    # Xóa các cột không phải là input của mô hình
    batch = {k: v for k, v in batch.items() if k in tokenizer.model_input_names}
    inputs = {k: torch.tensor(v).to(device) for k, v in batch.items()}
    with torch.no_grad():
        # Lấy hidden state của token [CLS]
        last_hidden_state = feature_extractor(**inputs).last_hidden_state
        cls_embedding = last_hidden_state[:, 0, :].cpu().numpy()
    return {"embeddings": cls_embedding}

embedded_dataset = sample_dataset.map(get_embeddings, batched=True, batch_size=16)

bert_embeddings = np.array(embedded_dataset['embeddings'])
labels = np.array(embedded_dataset['label'])

# Áp dụng t-SNE
from sklearn.manifold import TSNE

# Giảm chiều bằng t-SNE
tsne = TSNE(n_components=2, random_state=RANDOM_STATE, perplexity=30, n_iter_without_progress=1000)
embeddings_2d = tsne.fit_transform(bert_embeddings)

# Trực quan hóa
plt.figure(figsize=(10, 8))
palette = sns.color_palette("bright", 2)
sns.scatterplot(x=embeddings_2d[:, 0], y=embeddings_2d[:, 1], hue=labels, legend='full', palette=palette)
plt.title('t-SNE visualization of BERT Embeddings')
plt.show()
\end{minted}

% =======================================================================
% RESOURCE
% =======================================================================
\subsection*{A.9 Resource}
\label{subsec:resources}

\paragraph{Link video thuyết trình (Youtube):} https://youtu.be/sOHMstidMQg

\paragraph{Link source code (GitHub):} https://github.com/khoaoe/HCMUS-MathAI-Final-BERT

\paragraph{Sơ đồ của repo GitHub}
\begin{minted}{console}
HCMUS-MathAI-Final-BERT/
    ├── code/
    │   ├── bert.py                 # From-scratch BERT implementation
    │   ├── bert_application.ipynb  # experiment notebook
    │   └── requirements.txt        # Dependencies
    ├── report/                    # LaTeX report files' components
    │   ├── content.tex             # main content
    │   ├── appendix.tex            
    │   ├── cite.tex
    │   └── ref.bib
    ├── presentation/               # LaTeX file for slides
    │   └── presentation.tex
    └── README.md
\end{minted}